\section{Theory}

Automata are mathematical models of computation. They are used in several areas of computer science and linguistics. Finite state automata, one of the most basic models, are used in, among other things, text processing and compilers.

\subsection{Formal Languages}
Automata are concerned with \emph{formal languages}. A language in our setting is nothing more than a set of strings, or \emph{words}, constructed from a given (finite) \emph{alphabet}. An alphabet is typically denoted by the Greek capital $\Sigma$ and consists of \emph{symbols} or characters. A string over an alphabet is nothing more than a finite sequence of symbols. If $A$ is a set, then the set of all finite sequences over $A$ is denoted $A^*$. Hence $A^*$ is the full language over alphabet $A$. 

One special symbol that is a member of $A^*$ for any $A$ is $\epsilon$. This character denotes the \emph{empty string} and is just the word of length 0.

An example of an alphabet is $\Sigma = \{0,1\}$ and one of the many possible languages over $\Sigma$ is $\mathcal{L}= \{0w\ |\ w \in \Sigma^*\}$. This language contains all binary strings starting with 0 and nothing more.

\subsection{Combining Languages}
We define several operations over languages to construct new ones.

\begin{definition}
	content...
\end{definition}
